\section{Notation and the objectives we must keep separate}
\begin{longtable}{p{0.21\linewidth}p{0.70\linewidth}}
\toprule Symbol & Meaning \\
\midrule\endhead
$q$ & Problem or prompt drawn from a task distribution \\
$s_t$, $h_t$ & Reasoning step and prefix $h_t=(q,s_1,\ldots,s_t)$ \\
$\tau$, $T$ & Complete trajectory and number of reasoning steps \\
$\pi_\theta$, $\pi_0$ & Policy being optimized and fixed reference policy \\
$z_t$ & State before action $a_t$; in text, the prompt and current prefix \\
$c_t$, $v_t$ & Local step-correctness and cumulative prefix-validity labels \\
$p_{\pi_c}(h_t)$ & Probability a fixed completer $\pi_c$ finishes correctly \\
$\hat r_\phi$, $r^*$ & Learned process reward and specified target process reward \\
$G$ & Explicit step-to-trajectory aggregation rule \\
$\hat R$, $R^*$ & Proxy trajectory reward and intended task utility \\
$\U$, $L$, $U$ & Reward uncertainty set and its lower and upper envelopes \\
$b$, $\alpha$, $\delta$ & Width; prediction miscoverage level; confidence-set failure probability \\
$\rho$, $\beta$ & Ambiguity-set radius and policy KL regularization coefficient \\
\bottomrule
\end{longtable}

\textbf{Critical distinction.} In general,
\[
 c_t\ne v_t\ne p_{\pi_c}(h_t),\qquad
 G(r^*(h_1),\ldots,r^*(h_T))\ne R^*(\tau).
\]
A valid algebraic manipulation can follow an earlier false assumption; a wrong
step can later be repaired; a correct prefix can be difficult for a weak model
to complete. Even perfectly estimating step correctness does not automatically
make an arbitrary aggregation equal final-answer accuracy.

\section{Reinforcement learning in the reasoning setting}
\subsection{MDPs, states, actions, and policies}
A finite-horizon Markov decision process specifies states, actions, transition
probabilities, rewards, an initial-state distribution, and a horizon. The Markov
property says that, conditional on the current state and action, the next-state
distribution does not additionally depend on earlier history:
\[
 \Prob(z_{t+1}\mid z_{0:t},a_{0:t})
 =P(z_{t+1}\mid z_t,a_t).
\]
It does not mean only the immediately preceding observation matters: the state
can encode the entire history. Text generation is Markov when the state includes
the full prompt and generated prefix. An action can be one token, one reasoning
step, or a tool call. Appending text is deterministic given the action; sampling
actions from the policy is stochastic. Tool environments may also have stochastic
or hidden state. A partially observable formulation is needed when observations
do not determine the environment state.

A stochastic policy maps a state to an action distribution. A deterministic
policy picks one action. A stationary policy depends on the current state, and
a nonstationary one may also depend on time. Time can be appended to the state
in a finite-horizon problem. Existence of stationary deterministic optima requires
the usual MDP conditions and objective; do not assert it for arbitrary partial
observability or nonlinear robust objectives.

\subsection{Returns, value, and credit assignment}
Use the convention that $r_t$ is earned by action $a_t$:
\[
 G_t^{\mathrm{sum}}=\sum_{k=t}^{T-1}\gamma^{k-t}r_k,
 \quad V^\pi(z)=\E_\pi[G_t^{\mathrm{sum}}\mid z_t=z],
 \quad Q^\pi(z,a)=\E_\pi[G_t^{\mathrm{sum}}\mid z_t=z,a_t=a].
\]
The advantage $A^\pi(z,a)=Q^\pi(z,a)-V^\pi(z)$ measures whether an action is
better than the policy's average action. For additive returns,
\[
 Q^\pi(z,a)=\E[r_t+\gamma V^\pi(z_{t+1})\mid z_t=z,a_t=a].
\]
Sparse terminal rewards make it hard to determine which earlier decisions
caused success. Dense process feedback may help, but repeated erroneous rewards
also accumulate. The classical Monte Carlo, temporal-difference, value-learning,
and actor--critic vocabulary is useful background; these notes use a finite
episodic reasoning model rather than reproduce all tabular algorithms. The
textbook is a background reference \cite{sutton}.

\subsection{Policy gradients, PPO, and GRPO}
For a fixed trajectory reward, the likelihood-ratio identity gives
\[
 \nabla_\theta\E_{\tau\sim\pi_\theta}[R(\tau)]
 =\E\left[R(\tau)\sum_t\nabla_\theta\log\pi_\theta(a_t\mid z_t)\right].
\]
Baselines reduce variance; return-to-go or advantage estimates improve credit
assignment. A critic estimates value; it is a separate source of approximation
error from the reward model.

PPO commonly uses the clipped surrogate \cite{ppo}
\[
 \E_t\left[\min\{\omega_t(\theta)\hat A_t,
 \clip(\omega_t(\theta),1-\epsilon,1+\epsilon)\hat A_t\}\right],
 \quad \omega_t(\theta)=\frac{\pi_\theta(a_t\mid z_t)}{\pi_{\rm old}(a_t\mid z_t)}.
\]
Clipping limits this surrogate's incentive for large individual updates; it is
not a hard guarantee that policy KL remains small. The policy can still move
far from its reference through many small updates.

GRPO estimates relative advantages from groups of sampled answers, avoiding a
separately learned value critic in its common formulation \cite{deepseekmath}.
A simple outcome-level version uses
\[
 \hat A_j=\frac{R_j-\bar R}{\operatorname{std}(R_1,\ldots,R_m)+\varepsilon}.
\]
Dense process rewards require a stated token/step assignment; broadcasting one
trajectory advantage to every token is not stepwise process credit assignment.
Group normalization can obscure raw reward scales and gives little signal when
all group rewards coincide. Document the actual variant in experiments.

\subsection{RLHF, RLVR, and inference-time search}
RLHF commonly trains a preference reward model and optimizes against it. RLVR
uses automatically checked outcomes, such as answers or test results. The
checker is only as reliable as its specification and implementation. PRMs can
provide dense supervision in either setting. Best-of-$N$ changes which sampled
trajectory is selected while keeping generator weights fixed:
\[
 \tau_{\rm sel}=\arg\max_{j\le N}\hat R(\tau_j),\quad \tau_j\sim\pi_0.
\]
Beam search additionally selects prefixes and changes the effective exploration
distribution. RL changes generator parameters, allowing repeated adaptation to
the reward's errors. Results in one regime need not transfer to another.

\section{What a process reward model does}
A discriminative PRM predicts correctness scores at step boundaries from the
question and preceding reasoning. A generative PRM emits critiques or labels.
Conditioning on preceding steps is common; ``stepwise'' does not mean the model
ignores all previous context. The reward head may output a logit, a probability,
or a free scalar. The domain and scale determine admissible uncertainty sets.

\subsection{Labels and training}
For binary labels, a standard loss is
\[
 \ell_{\rm BCE}(\phi;h_t,y_t)
 =-y_t\log\hat r_\phi(h_t)-(1-y_t)\log(1-\hat r_\phi(h_t)).
\]
Human annotation may stop at the first error. Later steps then have missing or
censored labels; treating them all as independently labeled local errors changes
the target. Automated judges introduce correlated annotation mistakes.

For $K$ independent completions under a fixed completer $\pi_c$,
\[
 \tilde p(h_t)=\frac1K\sum_{k=1}^K\ind\{\text{completion }k\text{ succeeds}\},
 \quad \Var(\tilde p\mid h_t)=\frac{p_{\pi_c}(1-p_{\pi_c})}{K}.
\]
A label for ``at least one success'' has mean $1-(1-p_{\pi_c})^K$, not
$p_{\pi_c}$. It changes with $K$ and can approach one for any nonzero success
probability. Model repair makes completion success different from correctness
of the present step. Thus Monte Carlo uncertainty and semantic label mismatch
are distinct problems.

\subsection{Aggregation creates additional incentives}
For probability-like scores $r_t\in[0,1]$, common choices are
\[
 G_{\rm sum}=\sum_t r_t,\quad G_{\rm mean}=T^{-1}\sum_t r_t,
 \quad G_{\rm min}=\min_t r_t,\quad G_{\rm prod}=\prod_t r_t.
\]
Sum rewards repetition of easy steps. Mean can dilute a bad step with easy ones.
Minimum emphasizes the weakest step but can reward short or vacuous responses
and rejects repair unless repair is explicitly modeled. Product penalizes length:
$0.95^{10}\approx0.60$ and $0.95^{30}\approx0.21$. It is a joint-correctness
probability only under an appropriate probabilistic interpretation; multiplying
ordinary marginal step scores does not justify an independence claim.

For bounded per-step perturbations $|r_t-\tilde r_t|\le\epsilon$, sum errors can
grow as $T\epsilon$, while minimum errors are bounded by $\epsilon$. This
deterministic observation concerns the aggregation target, not proof that
minimum rewards optimize true mathematical accuracy. Splitting or merging steps
can change all four scores. Include segmentation attacks and empty-output rules.

\section{Why reward hacking occurs and why it matters}
Write $\hat R(\tau)=R^*(\tau)+e(\tau)$ schematically. Maximizing $\hat R$
selects both genuine quality and positive error. Error on ordinary data is not
the same as error among the candidates chosen by the optimizer.

\subsection{A numerical example}
Suppose a correct solution scores $0.80$ while an incorrect fluent solution
scores $0.95$. A policy increasing the latter's probability improves proxy
reward and reduces correctness. A useful tube might give the correct solution
$[0.76,0.84]$ and the exploitable solution $[0.20,1.00]$, making pessimistic
selection prefer the correct solution. If both intervals are too narrow, or the
wrong solution's uncertainty is underestimated, this defense fails. If every
interval is $[0,1]$, the lower score gives no useful ranking.

\subsection{Different mechanisms need different diagnostics}
\begin{itemize}
\item \textbf{Selection of noise:} higher search budgets find larger positive
errors even without new policy weights.
\item \textbf{Distribution shift:} optimization visits reasoning prefixes absent
from training, where feature--correctness associations fail.
\item \textbf{Misspecification:} a scorer rewards length, style, or easy local
claims instead of the intended reasoning property.
\item \textbf{Objective mismatch:} perfectly estimated local rewards are summed
or transformed into a trajectory objective that rewards the wrong behavior.
\item \textbf{Evaluator manipulation:} delimiters, injected instructions, or
checker loopholes change what the evaluator measures.
\end{itemize}
Optimization does not require conscious intent to exploit a reward. Conversely,
an ordinary wrong answer is not by itself evidence of hacking. Look for reward
improvement coupled to independent quality stagnation or decline and a mechanism
that explains it. Hacking harms model quality, invalidates capability estimates,
wastes compute, and can train persistent undesirable behavior.

\section{Uncertainty, calibration, and coverage}
\subsection{Three kinds of uncertainty}
\textbf{Epistemic uncertainty} concerns limited knowledge of the target reward
function. \textbf{Aleatoric variability} concerns stochastic outcomes or ambiguous
labels given the input. \textbf{Misspecification} means the model class or label
definition systematically misses relevant truth. Ensemble variance may detect
the first while missing the third. Rollout variance primarily measures outcome
variability and finite-sample noise, not intrinsic logical ambiguity.

\subsection{Prediction intervals are not confidence bands for a function}
An interval covering a new binary label is a prediction set for $Y$. A confidence
interval for $p(x)=\Prob(Y=1\mid X=x)$ covers a conditional mean. A function
confidence set asserts simultaneous containment of $p$ over a domain. They
require different information and assumptions. A binary-label interval that
contains both zero and one is valid but gives little reward discrimination.
Quantiles of observed labels or rollout estimates do not automatically become
confidence bounds on intrinsic correctness probabilities.

Calibration of a probability score $\mu$ means that among examples assigned a
score near $u$, the event frequency is near $u$. This is distribution-dependent.
Low aggregate expected calibration error can coexist with poor accuracy on
rare high-scoring mistakes. Ranking accuracy and calibration are also different.

Marginal coverage is
\[
 \Prob\{Y_{\rm new}\in C(X_{\rm new})\}\ge1-\alpha.
\]
Trajectory coverage contains all relevant steps. Candidate-set coverage contains
all trajectories being searched. Selection coverage concerns the chosen one.
Function-level coverage contains the target reward on every point of a stated
domain. Repeated RL adaptation requires an event valid for all encountered
queries or a sequentially justified update procedure; a static marginal claim
does not provide it.

\section{Robust optimization: the geometry matters}
For $\hat r:\mathcal H\to[0,1]$, a pointwise box tube is
\[
 \U_{\rm box}=\{r: |r(h)-\hat r(h)|\le b(h),\ h\in\mathcal H\}.
\]
It allows each point's error to vary independently. A function set induced by
one uncertain parameter $w$ couples different prefixes. An adversary choosing
one reward function for the entire policy gives
\[
 \max_\pi\inf_{r\in\U}\E_{\tau\sim\pi}[G(r(h_1),\ldots,r(h_T))].
\]
An adversary choosing a different model after each trajectory gives
\[
 \max_\pi\E_{\tau\sim\pi}\left[\inf_{r\in\U}G(r(h_1),\ldots,r(h_T))\right].
\]
Generally the second is no larger than the first. They can coincide for certain
rectangular sets but need not for coupled sets. State explicitly which is solved.

Distributional robustness instead permits an adversary to change a distribution.
For robust training this might be
\[
 \min_\phi\sup_{Q\in\mathcal Q}\E_Q[\ell_{\rm PRM}(\phi;h,Y)].
\]
It changes the fitted model, whereas policy pessimism changes how a fixed model
is used. Statistical DRO over examples, DRO over uncertain reward values, and
robust control over functions should not share one undifferentiated symbol.

\section{Distribution movement and reward shaping}
For $\gamma<1$, define normalized discounted occupancy
\[
 d_\pi(z,a)=(1-\gamma)\E_\pi\sum_{t\ge0}\gamma^t
 \ind\{z_t=z,a_t=a\}.
\]
Additive expected return is $(1-\gamma)^{-1}\langle d_\pi,r\rangle$.
For finite-horizon reasoning with $\gamma=1$, use the finite sum of state-action
visitation measures instead. Policy action KL, trajectory KL, and occupancy
divergence are different quantities. Under shared dynamics, trajectory KL can
be expressed as the expected sum of conditional action KLs, but an empirically
averaged token KL still needs normalization and sampling conventions.

An error bound averaged under $d_{\pi_0}$ may be amplified under $d_\pi$.
Importance weighting requires $d_\pi\ll d_{\pi_0}$ and a usable density ratio;
theoretical support overlap can be practically useless if weights are enormous.

With a reliable base task reward, potential shaping uses
\[
 F_t=\gamma\Phi(z_{t+1})-\Phi(z_t).
\]
The discounted sum telescopes:
\[
 \sum_{t=0}^{T-1}\gamma^tF_t=-\Phi(z_0)+\gamma^T\Phi(z_T).
\]
Setting the terminal potential to zero removes dependence on chosen actions
\cite{shaping}. Fix the potential during each evaluated episode and handle
truncation explicitly. Arbitrary dense PRM bonuses, adaptive within-episode
potentials, or nonzero length-dependent terminal potentials do not inherit this
argument.

\section{Research habits that make the results interpretable}
Fix the target and label protocol before selecting an uncertainty estimator.
Split data by problem, retain independent calibration and evaluation sets, and
keep attack construction separate from final transfer tests. Report correctness,
process validity, uncertainty coverage, computational cost, and reward curves.
An uncertain reward estimator cannot establish its own success merely by giving
higher scores. The detailed experiment plan and decision rules follow in Part III.
